\section{Experiments}
\vspace{-1ex}
\paragraph{Data.} We use the GEOM dataset \citep{GEOM}, consisting of QM9 (average 11 atoms) and DRUGS (average 44 atoms), to train and evaluate our model. 
We use the same train/val/test molecule splits, 106586/13323/1000 for
QM9 and 243473/30433/1000 for DRUGS, from \citet{GeoMol}. 

\vspace{-1ex}
\paragraph{Baselines.} We compare with Torsional Diffusion~\citep{torsional-diffusion}, which has outperformed all previous models on the previous evaluation metrics, and RDKit EKTDG + MMFF. We focus on these two methods as they depend on RDKit-generated structures directly or as input. 
For reference, Torsional Diffusion (TD) takes in the RDKit approximation and is trained to generate optimal torsion angle updates instead of explicit coordinates. 
We evaluate all methods over an expanded set of evaluation metrics to test the robustness of the generated conformations. 
For \model model configuration and compute resource breakdown, see Appendix \sref{section:appendix-breakdown}. 

\vspace{-1ex}
\paragraph{Data-constrained setting.} In real-world applications like polymer design, the availability of data is frequently limited, and accompanied by a scarcity of conformers for each molecule. The current datasets, QM9 and DRUGS, do not mimic this setting very well. For example, on average QM9 has 15 conformers per molecule and DRUGS has 104 per molecule---both datasets have significantly more conformers than in an experimental drug design setting. Given this, rather than training on the first 30 conformers as done in TD, we train on the first five and two conformers (typically those with the largest Boltzmann weight) for QM9 and DRUGS, respectively. This corresponds to 41\% of TD's training set or 23\% of the overall dataset for QM9, and 7.3\% of TD's training set or 1.7\% of DRUGS overall.
Even with less data, we outperform models trained on significantly larger training sets.

\begin{table}[t]
\caption{Quality of generated conformer ensembles for the GEOM-QM9 test set ($\delta=0.5$\AA) in terms of Coverage (\%) and Average RMSD (\AA). Torsional Diffusion (TD) was benchmarked using its evaluation code and available generated molecules, per their public instructions. Note that \model (5 epochs) was restricted to using 41\% of the data used by TD (250 epochs) to exemplify a low-compute and data-constrained setting. See Appendix \ref{section:appendix-qm9} for more details.}
\label{tab:results_qm9}
\begin{tabular}{lc|cccc|cccc} \toprule
               & \multicolumn{1}{l|}{} & \multicolumn{4}{c|}{Recall} & \multicolumn{4}{c}{Precision} \\
                 & \multicolumn{1}{l|}{} & \multicolumn{2}{c}{Coverage $\uparrow$} & \multicolumn{2}{c|}{AR $\downarrow$} & \multicolumn{2}{c}{Coverage $\uparrow$} & \multicolumn{2}{c}{AR $\downarrow$} \\
Method  & Metric & Mean & Med & Mean & Med & Mean & Med & Mean & Med \\ \midrule
\multicolumn{1}{l|}{RDKit + MMFF} & min  & 75.2 & \textbf{100.0} & 0.219 & 0.173 & \textbf{82.1} & \textbf{100.0} & \textbf{0.157} & \textbf{0.119}\\
% TD no filter  88.6 100 0.177 0.146 84.6 100 0.219 0.194
\multicolumn{1}{l|}{Torsional Diffusion} & min & \textbf{82.2} & \textbf{100.0} & \textbf{0.179} & \textbf{0.148} & 78.4 & \textbf{100.0} & 0.222  & 0.197\\
% \model (min) & 81.6 & \textbf{100.0} & 0.285 & 0.249 & 88.2 & \textbf{100.0} & 0.236 & 0.194 \\ no normalization filtering
\multicolumn{1}{l|}{\model} & min & 76.9 & \textbf{100.0} & 0.246 & 0.211 & 80.2 & \textbf{100.0} & 0.227 & 0.186\\
\midrule
% \midrule
\multicolumn{1}{l|}{RDKit + MMFF} & mean &  47.6 & \textbf{50.0} & 0.476 & 0.441 & 48.3 & 40.0 & 0.476 & 0.441 \\
\multicolumn{1}{l|}{Torsional Diffusion} & mean & 44.6 & 25.0 & 0.516 & 0.482 & 43.9 & 20.5 & 0.520 & 0.482\\
% \model (mean) & 54.8 & 64.6 & 0.471 & 0.433 & 55.2 & 92.3 & 0.471 & 0.433 \\
% \model (mean) & 48.8 & \textbf{50.0} & 0.466 & 0.434 & 49.7 & \textbf{50.0} & 0.466 & 0.434 \\ % 3 epoch
\multicolumn{1}{l|}{\model} & mean & \textbf{49.1} & \textbf{50.0} & \textbf{0.464} & \textbf{0.431} & \textbf{50.1} & \textbf{50.0} & \textbf{0.464} & \textbf{0.431}\\ % 5 epoch
% \midrule
\midrule
\multicolumn{1}{l|}{RDKit + MMFF} & max & 33.4 & 0.00 & 0.665 & 0.634 & 30.9 & 0.00 & 0.698 & 0.662\\
\multicolumn{1}{l|}{Torsional Diffusion} & max & 25.3 & 0.00 & 0.786 & 0.729 & 28.1 & 0.00 & 0.736 & 0.683\\
% \model (max) & 40.7 & 0.00 & 0.627 & 0.590 & 37.9 & 0.00 & 0.667 & 0.623 \\
% \model (max) & 35.5 & 0.00 & 0.639 & 0.604 & 33.8 & 0.00 & 0.647 & 0.625 \\ 3 epoch
\multicolumn{1}{l|}{\model} & max & \textbf{35.8} & 0.00 & \textbf{0.639} & \textbf{0.606} & \textbf{34.0} & 0.00 & \textbf{0.647} & \textbf{0.624}\\
% \midrule
\bottomrule
\end{tabular}
\end{table}

% \subsection{Metrics}
\vspace{-1ex}
\paragraph{Metrics.} We report both RMSD (AR) and Coverage for Recall (R) and Precision (P). Recall measures the number of ground truth conformers that are correctly predicted, and Precision measures the accuracy of the generated conformers.
Following \citet{torsional-diffusion}, we generate two times the number of ground truth conformers for each molecule. More formally, for $K=2L$ let $\{C^*_l\}_{l \in [1, L]}$ and $\{C_k\}_{k \in [1, K]}$ be respectively the sets of ground truth and generated conformers:
% \vspace{-4mm}
\begin{equation}
\begin{aligned}
% \text{COV-R} &:= \frac{1}{L}\, \bigg\lvert\{ l \in [1..L]: \exists k \in [1..K], \rmsd(C_k, C^*_l) < \delta \, \bigg\rvert\\
&\text{COV($f$)-R} := \frac{1}{L}\, \bigg\lvert\{ l \in [1..L]: f_{ k\in [1..K]} \rmsd(C_k, C^*_l) < \delta \} \bigg\rvert,\\
&\text{AR($f$)-R} := \frac{1}{L} \sum_{l \in [1..L]}  f_{ k\in [1..K]} \rmsd(C_k, C^*_l),\\
&\text{$f()$} := \{{\min, \mean, \max}\},
\end{aligned}
\end{equation}
% \vspace{-1mm}
where $\delta$ is the coverage threshold. The precision metrics are obtained by swapping ground truth and generated conformers. 
During a true test-time evaluation, the ground truth conformation is unknown. Therefore, we cannot compare it to all $2L$ generated structures to determine the most accurate one for the downstream application.
Thus, the additions of average (mean) and worst-case (max) evaluation allow us to measure the robustness of the generative model by providing an error bound on each sample. 
As most downstream tasks that use molecular conformers possess expensive oracle functions, \ie\ protein docking, it is important to be able to generate optimal conformers, as testing all $2L$ samples for each molecule may be infeasible.
We note that AR is the central task-agnostic metric. Coverage is derived from AR and its threshold can be chosen arbitrarily, with its value being task-dependent.

\vspace{-2ex}
\subsection{Results}
\label{section:results}
\vspace{-1ex}

\paragraph{QM9.} Our results, shown in~\tref{tab:results_qm9}, demonstrate a significant improvement over both RDKit+MMFF and Torsional Diffusion (TD) on the mean and max metrics.
We have effectively shown that our model can improve upon RDKit initializations, resulting in state-of-the-art generalization. 
\model does so using 59\% of TD's data and for only 5 epochs compared to 250.

\begin{table}[t]
\caption{Quality of generated conformer ensembles for the GEOM-DRUGS test set ($\delta=0.75$\AA) in terms of Coverage (\%) and Average RMSD (\AA). \model (5 epochs) was restricted to using 7.3\% of the data used by Torsional Diffusion (250 epochs) to exemplify a low-compute and data-constrained regime. We also benchmark GeoDiff using its generated molecules from \citet{torsional-diffusion} and GeoMol using its published weights \citep{GeoMol}. Further results in Appendix \sref{section:appendix-more-drugs}.} \label{tab:results_drugs}
\begin{tabular}{lc|cccc|cccc} \toprule
               & \multicolumn{1}{l|}{} & \multicolumn{4}{c|}{Recall} & \multicolumn{4}{c}{Precision}  \\
                 & \multicolumn{1}{l|}{} & \multicolumn{2}{c}{Coverage $\uparrow$} & \multicolumn{2}{c|}{AR $\downarrow$} & \multicolumn{2}{c}{Coverage $\uparrow$} & \multicolumn{2}{c}{AR $\downarrow$} \\
Method  & Metric & Mean & Med & Mean & Med & Mean & Med & Mean & Med \\ \midrule
\multicolumn{1}{l|}{RDKit + MMFF} & min  & 49.0 & 46.2 & 0.896 & 0.825 & 55.3 & 55.1 & 0.840 & \textbf{0.715} \\
\multicolumn{1}{l|}{GeoMol} & min  & 45.3 & 42.8 & 0.866 & 0.826 & 43.4 & 42.8 & 0.919 & 0.842 \\
\multicolumn{1}{l|}{GeoDiff} & min  & 43.1 & 39.3 & 0.828 & 0.806 & 25.5 & 15.2 & 1.131 & 1.083 \\
\multicolumn{1}{l|}{Torsional Diffusion} & min & \textbf{73.5} & \textbf{80.6} & \textbf{0.574} & \textbf{0.552} & \textbf{55.9} & \textbf{57.6} & \textbf{0.770} & 0.720 \\
% \model (min) & 81.6 & \textbf{100.0} & 0.285 & 0.249 & 88.2 & \textbf{100.0} & 0.236 & 0.194 \\ no normalization filtering
% \multicolumn{1}{l|}{\model full early} & min & 16.3 & 0.0 & 1.148 & 1.128 & 14.3 & 0.0 & 1.252 & 1.211 \\
\multicolumn{1}{l|}{\model} & min & 45.7 & 42.4 & 0.888 & 0.848 & 47.0 & 41.8 & 0.915 & 0.823 \\
\midrule
% \midrule
\multicolumn{1}{l|}{RDKit + MMFF} & mean & \textbf{8.87} & 0.00 & 1.903 & 1.887 & \textbf{8.34} & 0.00 & 1.903 & 1.887 \\
\multicolumn{1}{l|}{GeoMol} & mean  & 7.27 & 0.00 & 1.847 & 1.892 & 6.89 & 0.00 & 1.847 & 1.892 \\
\multicolumn{1}{l|}{GeoDiff} & mean  & 3.64 & 0.00 & 1.985 & 2.033 & 4.02 & 0.00 & 1.985 & 2.033 \\
\multicolumn{1}{l|}{Torsional Diffusion} & mean & 7.19 & 0.00 & 1.793 & 1.864 & 7.90 & 0.00 & 1.793 & 1.864\\
% \multicolumn{1}{l|}{\model full early} & mean & 5.01 & 0.0 & 1.777 & 1.789 & 4.74 & 0.0 & 1.777 & 1.789 \\
\multicolumn{1}{l|}{\model} & mean & 8.51 & 0.00 & \textbf{1.791} & \textbf{1.825} & 7.83 & 0.00 & \textbf{1.791} & \textbf{1.825} \\
% \midrule
\midrule
\multicolumn{1}{l|}{RDKit + MMFF} & max & \textbf{4.16} & 0.00 & 2.589 & 2.551 & 3.95 & 0.00 & 2.734 & 2.774 \\
\multicolumn{1}{l|}{GeoMol} & max  & 3.23 & 0.00 & 2.739 & 2.804 & 3.73 & 0.00 & 2.633 & 2.727 \\
\multicolumn{1}{l|}{GeoDiff} & max  & 1.66 & 0.00 & 2.946 & 3.062 & 2.28 & 0.00 & 2.711 & 2.806 \\
\multicolumn{1}{l|}{Torsional Diffusion} & max & 3.44 & 0.00 & 2.813 & 2.930 & \textbf{3.96} & 0.00 & 2.615 & 2.726 \\
% \multicolumn{1}{l|}{\model early} & max & 2.51 & 0.0 & 2.400 & 2.414 & 2.86 & 0.0 & 2.348 & 2.407 \\
\multicolumn{1}{l|}{\model} & max & 3.97 & 0.00 & \textbf{2.481} & \textbf{2.505} & 3.85 & 0.00 & \textbf{2.521} & \textbf{2.602} \\
% \midrule
\bottomrule
\end{tabular}
\end{table}

\vspace{-1ex}
\paragraph{DRUGS.} \tref{tab:results_drugs} demonstrates \modelns's ability to best generate robust low-error conformations for large drug-like molecules. With even further constraints on the data (only two conformers per molecule) and total iterations (5 epochs), \model achieves superior AR results for the more comprehensive mean and max metrics. For comparison, TD uses 30 conformers for 250 epochs.
We also demonstrate that \model outperforms GeoDiff~\citep{geodiff}, a diffusion model over 3D coordinate point clouds, and GeoMol, which learns to predict a minimal set of geometric quantities ($\ie$  angles
and bond lengths) and then reconstruct the local and global structures of the conformation.

\vspace{-1ex}
\paragraph{Takeaways.}
Both \model and TD optimize $p(X|\gR)$ but utilize the RDKit approximations $\gR$ in different ways. TD learns to update the torsion angles of $\gR$ while \model learns which geometric updates (coordinates, distances, and torsion angles) are optimal to translate $\gR$ to $X$.
Unlike TD, which uses a preprocessing optimization procedure to generate substitute ground truth conformers that mimic $p(\gR)$, 
\model directly learns from both $X$ and $\gR$ through its hierarchical graph matching procedure. We thus directly address the distributional shift problem via our learning framework.
We hypothesize that this, along with our torsion angle-based CG strategy, leads to our observed improvements while only using 2\% of the training iterations and 92.7\% less data compared to TD in the case of GEOM-DRUGS.
Although initialized with RDKit, \model and TD are slightly worse on some of the min metrics. We hypothesize that by querying the force field optimization $2L$ times, eventually, RDKit generates a single accurate structure whose identity is only known in this idealistic evaluation setting where the ground truth is known. 
\modelns, by contrast, is trained with a severely constrained number of conformations to generate $2L$ low-error structures over the entire ensemble.
We also note that min performance is not correlated with robust results in the mean and max metrics.
Overall, \model provides a comprehensive learning framework for robust conformer generation. 
Additional QM9 and DRUGS results are in Appendix \sref{section:appendix-qm9} and \sref{section:appendix-more-drugs}.
\vspace{-2ex}